\section*{Part IV: Break-Even, Runway, Reserves, and the Choice of Commitments}

\subsection*{33. Break-Even Is an Equation About a Specific Cost Structure}

In a simple single-offer model, let price be $P$, variable cost per unit $v$, fixed period costs $F$, and volume $Q$. Operating contribution after those fixed costs is:
\begin{equation}
\Pi = (P - v)Q - F.
\end{equation}
If $P > v$, the break-even volume is:
\begin{equation}
Q_{\text{BE}} = \left\lceil \frac{F}{P - v} \right\rceil.
\end{equation}

Suppose a technical cohort charges ₹30,000 per learner, incurs ₹8,000 of truly incremental acquisition and delivery costs per learner, and carries ₹2,20,000 of cohort-fixed delivery and administration cost. Contribution per learner is:
\begin{equation}
P - v = ₹30,000 - ₹8,000 = \mathbf{₹22,000}.
\end{equation}
Simple break-even is:
\begin{equation}
Q_{\text{BE}} = \frac{₹2,20,000}{₹22,000} = \mathbf{10\text{ learners}}.
\end{equation}
9 learners lose ₹22,000 on this model; 10 break even; 16 generate ₹1,32,000 before other enterprise costs.

But the calculation is only as good as the cost definition. If the founder's time is excluded, the answer describes \emph{cash contribution}, not normalized sustainability. If payment fees, refunds, or extra mentor labor rise with enrollment, $v$ is not constant. If the offering has multiple prices or scholarships, use realized price rather than the advertised tuition.

\subsection*{34. The Step-Cost Cliff Changes the Economics of Growth}

Extend the fictional cohort model: existing staff can responsibly support 16 learners. Above that number, the institute needs an additional ₹60,000 reviewer commitment for the cohort.
\begin{itemize}
\item At 16 learners, contribution after ₹2,20,000 fixed costs is:
\begin{equation}
\begin{aligned}
\Pi(16) &= (16 \times ₹22,000) - ₹2,20,000 \\
&= ₹3,52,000 - ₹2,20,000 = \mathbf{₹1,32,000}.
\end{aligned}
\end{equation}
\item At 17 learners, revenue net of per-learner variable costs rises to ₹3,74,000, but fixed costs jump to ₹2,80,000 (₹2,20,000 $+$ ₹60,000), leaving:
\begin{equation}
\begin{aligned}
\Pi(17) &= (17 \times ₹22,000) - ₹2,80,000 \\
&= ₹3,74,000 - ₹2,80,000 = \mathbf{₹94,000}.
\end{aligned}
\end{equation}
\end{itemize}

The 17th sale increased revenue and reduced profit by \textbf{₹38,000} in this discrete example. The next few learners may restore profitability as the new capacity fills, but the step is real. This is a reason to model capacity bands, not a reason to deny a learner necessary support.

A founder who promises unlimited enrollment based on average costs will eventually encounter these cliffs. Financial survival requires anticipating the \textbf{cash cost and hiring lead time} of the next capacity step before it becomes unavoidable.

\subsection*{35. Operating Break-Even, Cash Break-Even, and Debt-Service Break-Even}

These are distinct thresholds:
\begin{enumerate}
\item \textbf{Accounting or operating break-even} occurs when recognized revenues cover the expenses included in the chosen measure.
\item \textbf{Cash operating break-even} occurs when operating receipts cover necessary operating disbursements over the measured period.
\item \textbf{Full cash break-even} includes debt principal repayments, capital replacement, tax, and other investments.
\item \textbf{Normalized economic break-even} adds the opportunity or replacement cost of donated founder work.
\end{enumerate}

A prepaid program may have cash receipts well ahead of service and appear cash-positive for months without having earned the revenue. A lender may require principal repayment even in a month reporting accounting profit. A business can report positive EBITDA and still have negative free cash flow because of equipment purchases or working-capital growth.

Whenever a team says ``we are break-even,'' ask \emph{which definition}, over \emph{which period}, with \emph{which payments excluded}, and under \emph{which delivery obligations}. The answer should be 1 sentence and a reconciled schedule, not a slogan.

\subsection*{36. Margin of Safety Is More Useful Than the Exact Break-Even Point}

If break-even requires 10 paid learners, a target of exactly 10 is fragile. 1 refund, a weak applicant, or an unexpected instructional expense can erase contribution. A founder should plan for a \textbf{margin of safety} in both units and cash: enrollment above the minimum, costs below the maximum acceptable level, and a reserve for rework and lower collections.

In cost-volume analysis:
\begin{equation}
\text{Margin of Safety} = Q_{\text{planned}} - Q_{\text{BE}}.
\end{equation}
At 16 planned enrollments and 10 break-even enrollments, the margin is 6 learners or 37.5 percent of planned volume. But it is not a bankruptcy probability and it can overstate comfort when a late staffing step or refund event is likely.

Calculate downside volume as a scenario rather than as a vague mood. If 10 of the 16 sales disappear, if student mix changes, or if a major fixed commitment cannot be canceled, the model should show the actual cash consequence.

\subsection*{37. Runway: A Useful Ratio, and a Dangerous Simplification}

A common approximation is:
\begin{equation}
\text{Runway} = \frac{\text{Usable Cash}}{\text{Monthly Net Burn}},
\end{equation}
provided burn is positive and reasonably stable. If a venture has ₹7,20,000 of accessible cash but decides it must preserve a minimum reserve of ₹2,40,000, its \emph{discretionary spending buffer} is:
\begin{equation}
₹7,20,000 - ₹2,40,000 = \mathbf{₹4,80,000}.
\end{equation}
At ₹80,000 monthly net burn, runway is:
\begin{equation}
\frac{₹4,80,000}{₹80,000} = \mathbf{6\text{ months}}.
\end{equation}
If burn rises to ₹1,20,000, it falls to 4 months.

This example assumes no additional collections, financing, step-costs, or liabilities beyond the stated burn. In reality, money moves unevenly. A large payable due next week can shorten practical runway sharply. A prepaid teaching business may temporarily show positive cash flow while its remaining delivery obligations consume the apparent cushion.

Treat this ratio as a quick temperature reading. The \textbf{cash forecast's 1st date below the minimum allowable balance} is usually a more actionable survival horizon than the neat quotient. If current burn is zero or cash flow is positive, ``infinite runway'' can be meaningless when debt maturities, tax obligations, and future program costs are looming.

\subsection*{38. Gross Burn and Net Burn Reveal Different Risks}

\textbf{Gross burn} refers to cash operating outflows before offsetting operating receipts, under a stated definition. \textbf{Net burn} is operating cash outflows less operating cash inflows, often including adjustments for recurring capital spending depending on convention. Define both and be consistent.

Suppose the company spends ₹2 lakh monthly and collects ₹1.4 lakh: gross burn is ₹2 lakh and net burn is ₹60,000 under this simplified operating convention. If a major payer fails to collect, gross burn remains and net burn could suddenly become ₹2 lakh. The distinction reveals dependence on continuing receipts.

A founder should track both the \textbf{expense commitments that must be paid regardless of sales} and the \textbf{actual collection performance}. A business that says its burn is small only because it expects a large unsigned contract is not being conservative. 1-off financing cash should not be counted as revenue that reduces operating burn.
